\chapter{Divergences}
\label{ch:info}

\section{Kullback--Leibler divergence}

\begin{defbox}
For densities $p$ and $q$ with $q>0$ wherever $p>0$,
\[
\KL(p\Vert q)=\E_{x\sim p}\Bigl[\log\frac{p(x)}{q(x)}\Bigr]
=\int p(x)\log\frac{p(x)}{q(x)}\,dx.
\]
\end{defbox}

\begin{resultbox}[Properties from Lecture 2, slides 36--39]
\begin{enumerate}[label=\arabic*.]
\item $\KL(p\Vert q)\ge 0$.
\item $\KL(p\Vert q)=0$ if and only if $p=q$ almost everywhere.
\item $\KL(p\Vert q)\neq \KL(q\Vert p)$ in general.
\item KL is not a distance: it fails symmetry and the triangle inequality.
\item KL may be infinite when the supports disagree.
\end{enumerate}
\end{resultbox}

The proof of nonnegativity is Jensen. Write
\[
-\KL(p\Vert q)=\E_p\bigl[\log(q/p)\bigr]\le \log \E_p[q/p]=\log\int q=0,
\]
so the KL is at least zero. Equality in Jensen holds when $q/p$ is constant on the support of $p$. Normalization forces that constant to be $1$.

\begin{pitbox}[Which way the KL points]
In a VAE the bound contains $\KL(q_\phi(z\given x)\Vert p_\theta(z))$. The first argument is the encoder, the second is the prior. Pushing that KL to zero makes the encoder match the prior, which is not the same as making the prior match the aggregated encoder. The asymmetry is the source of posterior collapse and of the off-mode mass of a Gaussian prior.
\end{pitbox}

\section{Closed-form Gaussian KL}

Let $q=\N(\mu_q,\diag(\sigma_q^2))$ and $p=\N(\mu_p,\diag(\sigma_p^2))$ on $\R^d$, with every $\sigma>0$. Independence across coordinates splits the KL into a sum. In one dimension, the Gaussian density contributes
\[
\log q(z)=-\tfrac12\log(2\pi\sigma_q^2)-\frac{(z-\mu_q)^2}{2\sigma_q^2}.
\]
Taking the expectation under $q$ and subtracting the same expression for $p$ gives, after using $\E_q[(z-\mu_q)^2]=\sigma_q^2$ and $\E_q[(z-\mu_p)^2]=\sigma_q^2+(\mu_q-\mu_p)^2$,
\begin{resultbox}[Diagonal Gaussian KL]
\begin{equation}
\label{eq:diagkl}
\KL(q\Vert p)=\frac12\sum_{j=1}^{d}
\Biggl(
\log\frac{\sigma_{p,j}^2}{\sigma_{q,j}^2}
-1
+\frac{\sigma_{q,j}^2+(\mu_{q,j}-\mu_{p,j})^2}{\sigma_{p,j}^2}
\Biggr).
\end{equation}
Against the standard normal $p=\N(0,I)$ this collapses to
\begin{equation}
\label{eq:stdkl}
\KL(q\Vert \N(0,I))=\frac12\sum_{j=1}^{d}\bigl(\mu_{q,j}^2+\sigma_{q,j}^2-1-\log\sigma_{q,j}^2\bigr).
\end{equation}
\end{resultbox}

Equation~\eqref{eq:stdkl} is the KL implemented in the VAE notebook. Equation~\eqref{eq:diagkl} is Homework~2, problem~1.3, before that specialization.

The same calculation with full covariances $\Sigma_0,\Sigma_1\succ 0$ produces the multivariate formula used in Homework~1, problem~1.4(b). Write $p_i=\N(\mu_i,\Sigma_i)$. Then
\begin{resultbox}[Multivariate Gaussian KL]
\label{res:mvnkl}
\begin{equation}
\label{eq:mvnkl}
\KL(p_0\Vert p_1)
=\frac12\Biggl[
\tr(\Sigma_1^{-1}\Sigma_0)
+(\mu_0-\mu_1)^\top\Sigma_1^{-1}(\mu_0-\mu_1)
-d
+\log\frac{\det\Sigma_1}{\det\Sigma_0}
\Biggr].
\end{equation}
\end{resultbox}
The trace and Mahalanobis terms come from $\E_{p_0}[(x-\mu_1)^\top\Sigma_1^{-1}(x-\mu_1)]$. Subtracting the differential entropy $h(p_0)=\tfrac d2\log(2\pi e)+\tfrac12\log\det\Sigma_0$ cancels the $\tfrac d2\log(2\pi)$ pieces and leaves~\eqref{eq:mvnkl}. Chapter~\ref{ch:gauss} uses this formula inside the maximum-entropy Pythagorean identity.

\begin{practbox}[A number you can check by hand]
Compute $\KL(\N(0,1)\Vert \N(0,e))$ and $\KL(\N(0,e)\Vert \N(0,1))$.
\end{practbox}
\begin{solbox}
Use $d=1$, $\mu=0$. First pair: $\sigma_q^2=1$, $\sigma_p^2=e$, so
$\tfrac12(\log e-1+1/e)=\tfrac12(1-1+e^{-1})=1/(2e)$.
Second pair: $\sigma_q^2=e$, $\sigma_p^2=1$, so
$\tfrac12(\log e^{-1}-1+e)=\tfrac12(-1-1+e)=(e-2)/2$.
They differ. KL is not symmetric, and the direction that puts a wide Gaussian on a narrow one is the more expensive one here.
\end{solbox}

\section{Jensen--Shannon divergence}

Let $m=\tfrac12(p+q)$. Lecture~2 defines
\[
\JS(p\Vert q)=\tfrac12\KL(p\Vert m)+\tfrac12\KL(q\Vert m).
\]
JS is symmetric. In nats, $0\le \JS(p\Vert q)\le \log 2$. The upper bound $1$ on the lecture slides is the same statement in bits, where the log is base $2$. In nats the sharp cap is $\ln 2$, not $1$. $\sqrt{\JS}$ is a metric. JS itself is not, because it can violate the triangle inequality.

\begin{intbox}
KL is infinite when supports are disjoint, so it cannot compare a model that puts no mass on a point with a data distribution that does. JS stays finite because $m$ charges both supports. GANs later exploit a divergence with that behaviour. Homework~1 asks you to see the failure of KL directly.
\end{intbox}

Expanding the two KL terms and rearranging Shannon entropies gives the entropy form
\begin{resultbox}[JS as an entropy gap]
\[
\JS(p\Vert q)=H(m)-\tfrac12 H(p)-\tfrac12 H(q),
\]
in the same base used for the logarithms. Equivalently, if $Z\sim\mathrm{Bernoulli}(1/2)$ and $X\mid Z=0\sim p$, $X\mid Z=1\sim q$, then $X\sim m$ and
\[
\JS(p\Vert q)=I(Z;X)=H(X)-H(X\given Z).
\]
\end{resultbox}
Mutual information satisfies $0\le I(Z;X)\le H(Z)=\log 2$ in nats. The lower bound is attained if and only if $p=q$. The upper bound is attained if and only if $Z$ is a function of $X$, i.e.\ $p$ and $q$ are mutually singular. Pointwise, $m\ge p/2$ on $\{p>0\}$ forces $\log(p/m)\le\log 2$, so each half of JS is at most $\log 2$ and the average recovers the same cap.

\begin{practbox}[JS of two point masses]
On $\{0,1\}$, take $p=\delta_0$ and $q=\delta_1$. Compute $\KL(p\Vert q)$ and $\JS(p\Vert q)$ in nats.
\end{practbox}
\begin{solbox}
$\KL(p\Vert q)=\infty$ because $q(0)=0$ while $p(0)=1$. The mixture $m$ is uniform on two atoms, $H(m)=\log 2$, and $H(p)=H(q)=0$, so $\JS=\log 2$. JS saturates while KL is infinite.
\end{solbox}

\section{Bernoulli KL near the boundary}

For Bernoulli parameters $a,c\in(0,1)$,
\[
\KL\bigl(\mathrm{Ber}(a)\Vert\mathrm{Ber}(c)\bigr)
=a\log\frac ac+(1-a)\log\frac{1-a}{1-c}.
\]
As $c\to 0^+$ with $a$ fixed in $(0,1)$, the expansion $-\log(1-c)=c+O(c^2)$ yields
\begin{resultbox}[Asymmetry at the edge]
\[
\KL\bigl(\mathrm{Ber}(a)\Vert\mathrm{Ber}(c)\bigr)=-a\log c-H(\mathrm{Ber}(a))+O(c)\to+\infty,
\]
while
\[
\KL\bigl(\mathrm{Ber}(c)\Vert\mathrm{Ber}(a)\bigr)\to-\log(1-a)\in(0,\infty).
\]
\end{resultbox}
Charging an atom that the other distribution abandons costs arbitrarily large KL in one direction and only a finite KL in the other. The contribution $c\log c$ of a vanishing atom disappears.

The same one-sided blow-up shows that KL fails the triangle inequality. For $1>a>b>0$ and $c\to 0^+$,
\[
\KL(\mathrm{Ber}(a)\Vert\mathrm{Ber}(c))-\KL(\mathrm{Ber}(b)\Vert\mathrm{Ber}(c))
\]
grows like $-(a-b)\log c\to+\infty$, while $\KL(\mathrm{Ber}(a)\Vert\mathrm{Ber}(b))$ stays finite. For every small enough $c$,
\[
\KL(\mathrm{Ber}(a)\Vert\mathrm{Ber}(c))
>\KL(\mathrm{Ber}(a)\Vert\mathrm{Ber}(b))+\KL(\mathrm{Ber}(b)\Vert\mathrm{Ber}(c)).
\]

\begin{practbox}[A numerical triangle break]
In nats, compare $\KL(\mathrm{Ber}(1/2)\Vert\mathrm{Ber}(1/100))$ with the sum through $\mathrm{Ber}(1/5)$.
\end{practbox}
\begin{solbox}
Direct evaluation gives about $1.6145$ on the left and $0.2231+0.4287=0.6518$ on the right. The direct path is more than twice as expensive as the broken path.
\end{solbox}

\section{Homework 1 problems}

The Monte Carlo, importance sampling, and perplexity parts of Homework~1 problem~2 live in Chapter~\ref{ch:mc}. The KL and JS parts are collected here. The multivariate Gaussian KL formula~\eqref{eq:mvnkl} is the identity invoked by Homework~1, problem~1.4(b); the full Pythagorean argument sits with the Gaussian maximum-entropy section in Chapter~\ref{ch:gauss}.

\begin{hwqbox}[Homework 1, problem 2.3]
The Jensen--Shannon divergence is $\JS(p\Vert q)=\tfrac12\KL(p\Vert m)+\tfrac12\KL(q\Vert m)$ with $m=(p+q)/2$. Rewrite JS in terms of Shannon entropies. Show that $\JS(p\Vert q)=I(Z;X)$ for a fair Bernoulli bit $Z$ that chooses between sampling $X$ from $p$ or from $q$. Deduce the sharp bounds $0\le\JS\le\log 2$ in nats, and characterize equality.
\end{hwqbox}
\begin{hwsolbox}[Homework 1, problem 2.3]
Expanding the two KL terms,
\begin{align*}
\JS(p\Vert q)
&=\tfrac12\sum p\log\frac pm+\tfrac12\sum q\log\frac qm\\
&=-\tfrac12 H(p)-\tfrac12 H(q)-\sum m\log m\\
&=H(m)-\tfrac12 H(p)-\tfrac12 H(q).
\end{align*}
Let $Z\sim\mathrm{Bernoulli}(1/2)$ and $X\mid Z=0\sim p$, $X\mid Z=1\sim q$. Then $X\sim m$ and
\[
I(Z;X)=H(X)-H(X\given Z)=H(m)-\tfrac12 H(p)-\tfrac12 H(q)=\JS(p\Vert q).
\]
In nats, $0\le I(Z;X)\le H(Z)=\log 2$. The lower bound holds with equality if and only if $Z\perp X$, i.e.\ $p=q$. The upper bound holds if and only if $H(Z\given X)=0$, i.e.\ $Z$ is a function of $X$, i.e.\ $p$ and $q$ are mutually singular.
\end{hwsolbox}

\begin{hwqbox}[Homework 1, problem 2.3(a)]
Prove the pointwise bound $\KL(p\Vert m)\le\log 2$ (and likewise for $q$), and recover $\JS\le\log 2$. Give an example where $\KL(p\Vert q)=\infty$ while $\JS(p\Vert q)$ is finite and maximal. Comment on the bound $0\le\JS\le 1$ that appears when the logarithm is taken in base~$2$.
\end{hwqbox}
\begin{hwsolbox}[Homework 1, problem 2.3(a)]
Since $m\ge p/2$, one has $\log(p/m)\le\log 2$ on $\{p>0\}$, so $\KL(p\Vert m)\le\log 2$. The same bound holds for $q$, and averaging recovers $\JS\le\log 2$. Equality holds simultaneously when $m=p/2$ $p$-a.s.\ and $m=q/2$ $q$-a.s., i.e.\ on disjoint supports.

If $p=\delta_0$ and $q=\delta_1$ on $\{0,1\}$, then $q(0)=0$ while $p(0)=1$, so $\KL(p\Vert q)=\infty$, but the supports are disjoint and $\JS(p\Vert q)=\log 2$, the maximum.

The quoted bound $0\le\JS\le 1$ is tight for the base-$2$ logarithm (the maximum is one bit). In nats the supremum is $\ln 2$, and the upper bound $1$ is valid but not sharp.
\end{hwsolbox}

\begin{hwqbox}[Homework 1, problem 2.4]
For $a\in(0,1)$ fixed, expand $\KL(\mathrm{Ber}(a)\Vert\mathrm{Ber}(c))$ as $c\to 0^+$, and evaluate the limit of $\KL(\mathrm{Ber}(c)\Vert\mathrm{Ber}(a))$. Explain the asymmetry.
\end{hwqbox}
\begin{hwsolbox}[Homework 1, problem 2.4]
\begin{align*}
\KL(\mathrm{Ber}(a)\Vert\mathrm{Ber}(c))
&=a\log\frac ac+(1-a)\log\frac{1-a}{1-c}\\
&=-a\log c+\bigl(a\log a+(1-a)\log(1-a)\bigr)-(1-a)\log(1-c).
\end{align*}
As $c\to 0^+$, $-\log(1-c)=c+O(c^2)$, so
\[
\KL(\mathrm{Ber}(a)\Vert\mathrm{Ber}(c))
=-a\log c-H(\mathrm{Ber}(a))+O(c).
\]
The leading term is $-a\log c\to+\infty$. The constant after it is $-H(\mathrm{Ber}(a))$.

In the other direction, $c\log c\to 0$ and $(1-c)\log((1-c)/(1-a))\to-\log(1-a)$, so
\[
\KL(\mathrm{Ber}(c)\Vert\mathrm{Ber}(a))\to-\log(1-a)\in(0,\infty).
\]
KL is not symmetric: charging the point that the other distribution is abandoning costs an arbitrarily large KL in one direction, and only a finite KL in the other. The $c\log c$ contribution of a vanishing atom disappears.
\end{hwsolbox}

\begin{hwqbox}[Homework 1, problem 2.4(a)]
Show that KL does not satisfy the triangle inequality, by exhibiting Bernoulli parameters $1>a>b>0$ and a sequence $c\to 0^+$ for which
\[
\KL(\mathrm{Ber}(a)\Vert\mathrm{Ber}(c))
>\KL(\mathrm{Ber}(a)\Vert\mathrm{Ber}(b))+\KL(\mathrm{Ber}(b)\Vert\mathrm{Ber}(c)).
\]
Give a concrete numerical triple in nats.
\end{hwqbox}
\begin{hwsolbox}[Homework 1, problem 2.4(a)]
For $1>a>b>0$,
\begin{align*}
\KL(\mathrm{Ber}(a)\Vert\mathrm{Ber}(c))&=-a\log c+O(1),\\
\KL(\mathrm{Ber}(b)\Vert\mathrm{Ber}(c))&=-b\log c+O(1),
\end{align*}
while $\KL(\mathrm{Ber}(a)\Vert\mathrm{Ber}(b))$ stays finite. Their difference grows like $-(a-b)\log c\to+\infty$. So for every sufficiently small $c>0$ the displayed strict inequality holds.

Numerical triple, in nats: $(a,b,c)=(1/2,\,1/5,\,1/100)$,
\[
\KL(\mathrm{Ber}(1/2)\Vert\mathrm{Ber}(1/100))\approx 1.6145,
\]
\[
\KL(\mathrm{Ber}(1/2)\Vert\mathrm{Ber}(1/5))+\KL(\mathrm{Ber}(1/5)\Vert\mathrm{Ber}(1/100))\approx 0.2231+0.4287=0.6518.
\]
\end{hwsolbox}
